%==============================================================
%  cir.tex
%  hasklib — mathematical companion / derivations
%
%  Cox-Ingersoll-Ross process: unlike gbm.tex/ou.tex/abm.tex, no
%  closed-form path solution is derived here. This file covers what
%  the code actually needs: moments (via an ODE method, not a path
%  solution), the Feller condition, and the motivation for the full
%  truncation numerical scheme.
%
%  Sibling of gbm.tex, ou.tex, abm.tex, not a dependent of any.
%==============================================================
\documentclass[11pt]{article}

%---------------- packages ----------------
\usepackage[utf8]{inputenc}
\usepackage[T1]{fontenc}
\usepackage{amsmath, amssymb, amsthm}
\usepackage{mathtools}
\usepackage[a4paper, margin=1in]{geometry}
\usepackage{xcolor}
\usepackage{hyperref}
\usepackage{cleveref}
\usepackage{bm}
\usepackage{harvard}   % Harvard author-date referencing (with agsm.bst)

%---------------- notation macros ----------------
% (kept identical to ito_foundations.tex / gbm.tex / ou.tex / abm.tex)
\newcommand{\Prob}{\mathbb{P}}
\newcommand{\Qmeas}{\mathbb{Q}}
\newcommand{\E}{\mathbb{E}}
\newcommand{\Var}{\operatorname{Var}}
\newcommand{\Cov}{\operatorname{Cov}}
\newcommand{\filt}{\mathcal{F}}
\newcommand{\Om}{\Omega}

\newcommand{\dd}{\mathrm{d}}
\newcommand{\dW}{\dd W_t}
\newcommand{\dt}{\dd t}
\newcommand{\R}{\mathbb{R}}

\newcommand{\drift}{a}
\newcommand{\diff}{b}

%---------------- theorem environments ----------------
\theoremstyle{plain}
\newtheorem{theorem}{Theorem}
\newtheorem{lemma}{Lemma}
\newtheorem{proposition}{Proposition}
\newtheorem{corollary}{Corollary}

\theoremstyle{definition}
\newtheorem{definition}{Definition}

\theoremstyle{remark}
\newtheorem{remark}{Remark}

%==============================================================
\title{The Cox--Ingersoll--Ross Process: Moments, the Feller\\
       Condition, and Full Truncation\\
       \large \texttt{hasklib} mathematical companion}
\author{Hubert}
\date{\today}
%==============================================================

\begin{document}
\maketitle

\begin{abstract}
Unlike GBM, OU, and ABM, the Cox--Ingersoll--Ross (CIR) process
\citeasnoun{cir1985} has no closed-form path solution derived in this
note: neither the $\ln x$ substitution nor an integrating factor
removes the $x$-dependence from its square-root diffusion coefficient.
What the library actually needs from this process --- its mean and
variance, the parameter condition (Feller's condition) guaranteeing it
stays strictly positive, and the motivation for the numerical scheme
used when that condition is in doubt --- is derived instead. The exact
transition law (a scaled noncentral chi-squared distribution) exists
but is not derived here; deriving it requires machinery (the
squared-Bessel-process connection, or a PDE/Laplace-transform
argument) well beyond this note's scope.
\end{abstract}

%--------------------------------------------------------------
\section{The CIR SDE}
%--------------------------------------------------------------

\begin{definition}[Cox--Ingersoll--Ross process]\label{def:cir}
$X_t$ follows a CIR process with mean-reversion speed $\kappa>0$,
long-run mean $\theta>0$, and volatility $\sigma>0$ if it solves
\begin{equation}\label{eq:cirsde}
   \dd X_t = \kappa(\theta - X_t)\,\dt + \sigma\sqrt{X_t}\,\dW,
   \qquad X_0>0.
\end{equation}
In the notation of \texttt{ito\_foundations.tex} this is
$\drift(t,x)=\kappa(\theta-x)$ (identical to OU) and
$\diff(t,x)=\sigma\sqrt{x}$.
\end{definition}

\begin{remark}[Correspondence with the code]
\texttt{CIR::drift} returns $\kappa(\theta-x)$, exactly matching
\texttt{OU::drift}. \texttt{CIR::diffusion} returns $\sigma\sqrt{x}$,
which requires $x\ge0$ to be well-defined --- unlike GBM's $\sigma x$
or OU's constant $\sigma$, this diffusion coefficient is not defined
for negative arguments at all, which is the source of every
complication in this file.
\end{remark}

\begin{remark}[A genuine gap in the existence--uniqueness story]
The existence--uniqueness assumption in \texttt{ito\_foundations.tex}
requires $\drift$ and $\diff$ to be Lipschitz. $\sqrt{x}$ is
\emph{not} Lipschitz near $x=0$ (its slope $\to\infty$ there), so
that assumption does not literally cover \eqref{eq:cirsde} as stated.
A strong solution nonetheless exists and is unique --- this requires a
separate, more delicate existence--uniqueness theory for
H\"older-continuous (rather than Lipschitz) diffusion coefficients;
see \citeasnoun{cir1985} or \citeasnoun{brigomercurio2006}. \emph{(Cited,
not proved, in the same spirit as the original assumption --- the
library depends only on the conclusion.)}
\end{remark}

%--------------------------------------------------------------
\section{Why the earlier tricks do not apply here}
%--------------------------------------------------------------

GBM's diffusion $\sigma x$ is proportional to the state, which
$\ln x$ removed. OU's diffusion $\sigma$ is constant, which the
integrating factor $e^{\kappa t}$ exploited alongside an affine drift.
CIR's drift $\kappa(\theta-x)$ is the same affine form as OU's, so the
integrating factor idea is tempting to reuse --- but it was the
diffusion term, not the drift, that made OU's trick work cleanly:
applying It\^o's lemma to $f(t,x)=e^{\kappa t}x$ needs $f_{xx}=0$ to
kill the second-order term entirely, which held because $f$ was
\emph{linear} in $x$. The moment we try to fold $\sigma\sqrt{x}$ into
any single substitution $f(t,x)$, the resulting $f_{xx}$ term no
longer vanishes, and no elementary choice of $f$ linearises
\eqref{eq:cirsde} the way $\ln x$ or $e^{\kappa t}x$ did for the
earlier processes. This is not a failure of effort --- $\sqrt{x}$ is a
genuinely different kind of nonlinearity (a \emph{square-root}, or
CEV-type, diffusion), and no algebraic trick of this kind exists for
it in general.

\begin{remark}[The exact law exists, but is out of scope here]
CIR's transition density is nonetheless known exactly: conditional on
$X_t$, the (scaled) law of $X_{t+\Delta t}$ is a noncentral
chi-squared distribution, a result following from CIR's connection to
squared Bessel processes \cite{cir1985}. This is the route
\texttt{CIR::sample\_terminal} would take for an exact sampler
(Boost.Math's noncentral chi-squared machinery, per the library's
\texttt{stats/} layer), analogous to \texttt{GBM::sample\_terminal}
and \texttt{OU::sample\_terminal}. Deriving that transition density,
however, requires either the squared-Bessel-process identification or
a PDE/Laplace-transform argument neither established nor needed
elsewhere in this note; it is cited, not derived, here.
\end{remark}

%--------------------------------------------------------------
\section{Moments, via an ODE rather than a path solution}
%--------------------------------------------------------------

Although no closed-form $X_t$ is derived, its first two moments are
obtainable directly: take expectations of the SDE (or of It\^o's
lemma applied to a power of $X_t$) and solve the resulting
\emph{deterministic} ODE. This sidesteps needing a path solution at
all.

\begin{proposition}[Mean of CIR]\label{prop:cir-mean}
Let $m(t)=\E[X_t\mid X_0]$. Then $m(t)$ solves
$m'(t)=\kappa(\theta-m(t))$ with $m(0)=X_0$, giving
\[
   m(t) = X_0\,e^{-\kappa t} + \theta\big(1-e^{-\kappa t}\big),
\]
identical in form to OU's mean (\texttt{ou.tex}).
\end{proposition}

\begin{proof}
Taking expectations of the integral form of \eqref{eq:cirsde},
\[
   m(t) = X_0 + \kappa\theta\int_0^t \dd s - \kappa\int_0^t m(s)\,\dd s
        + \E\left[\int_0^t \sigma\sqrt{X_s}\,\dd W_s\right].
\]
The stochastic integral has zero expectation (the adapted It\^o
isometry cited in \texttt{schemes.tex}; \citeasnoun{shreve2004},
\citeasnoun{oksendal2003}), regardless of the specific form of the
diffusion coefficient inside it. Differentiating the remaining
deterministic relation gives $m'(t)=\kappa\theta-\kappa m(t)$, exactly
the deterministic ODE behind OU's mean, with the same solution.
\end{proof}

\begin{remark}
The mean depends only on the drift, and CIR's drift is literally the
same function as OU's; the two processes' means therefore coincide
exactly, even though their diffusion coefficients --- and hence their
variances, paths, and positivity properties --- are entirely
different.
\end{remark}

\begin{proposition}[Variance of CIR]\label{prop:cir-var}
Let $v(t)=\Var(X_t\mid X_0)$. Then
\[
   v(t) = \frac{\sigma^2 X_0}{\kappa}\Big(e^{-\kappa t}-e^{-2\kappa t}\Big)
        + \frac{\sigma^2\theta}{2\kappa}\Big(1-e^{-\kappa t}\Big)^2 .
\]
\end{proposition}

\begin{proof}
Apply It\^o's lemma to $f(x)=x^2$: $f_x=2x$, $f_{xx}=2$, $f_t=0$.
With $\drift=\kappa(\theta-X_t)$ and $\diff=\sigma\sqrt{X_t}$,
\[
   \dd(X_t^2)
     = \Big(2X_t\cdot\kappa(\theta-X_t) + \tfrac12\cdot\sigma^2 X_t\cdot 2\Big)\dt
       + 2X_t\cdot\sigma\sqrt{X_t}\,\dW
     = \big(2\kappa\theta X_t - 2\kappa X_t^2 + \sigma^2 X_t\big)\dt
       + 2\sigma X_t^{3/2}\,\dW .
\]
Taking expectations and writing $w(t)=\E[X_t^2]$, the stochastic term
again vanishes in expectation, giving
\[
   w'(t) = \big(2\kappa\theta+\sigma^2\big)m(t) - 2\kappa\,w(t).
\]
Rather than solve for $w(t)$ directly, substitute $w(t)=v(t)+m(t)^2$
and use $m'(t)=\kappa(\theta-m(t))$ from \Cref{prop:cir-mean} to
simplify. Differentiating $w=v+m^2$ gives
$w'=v'+2mm'=v'+2\kappa\theta m-2\kappa m^2$. Substituting into the
equation for $w'(t)$ and cancelling the $2\kappa\theta\,m(t)$ terms on
both sides,
\[
   v'(t) + 2\kappa\theta\,m(t) - 2\kappa\,m(t)^2 + 2\kappa\big(v(t)+m(t)^2\big)
     = \big(2\kappa\theta+\sigma^2\big)m(t)
   \;\;\Longrightarrow\;\;
   v'(t) + 2\kappa\,v(t) = \sigma^2 m(t),
\]
with $v(0)=0$ (since $X_0$ is a fixed starting point, not itself
random). This is a much cleaner ODE than the one for $w(t)$: the
$\theta$-dependence has been absorbed entirely into $m(t)$, which is
already known. Solve by the same integrating-factor technique used
for OU's SDE, now applied to this ordinary (non-stochastic) ODE:
multiplying by $e^{2\kappa t}$ collapses the left side into
$\tfrac{\dd}{\dd t}\big(e^{2\kappa t}v(t)\big)$, and substituting
$m(t)$ from \Cref{prop:cir-mean} and integrating from $0$ to $t$
(using $v(0)=0$) gives, after simplification,
\[
   v(t) = \frac{\sigma^2 X_0}{\kappa}\big(e^{-\kappa t}-e^{-2\kappa t}\big)
        + \frac{\sigma^2\theta}{2\kappa}\big(1-e^{-\kappa t}\big)^2,
\]
which is the claimed formula. \emph{(The intermediate integrating-factor
algebra is routine and omitted; it is the same mechanical steps as the
OU variance derivation in \texttt{ou.tex}, applied to this ODE instead
of directly to the SDE.)}
\end{proof}

\begin{remark}[What this moment computation buys, and what it does not]
\Cref{prop:cir-mean,prop:cir-var} give exact first and second moments
with no approximation --- useful for moment-matching tests analogous
to \texttt{test\_stochastic.cpp}'s GBM check. They do \emph{not} give
the full distribution of $X_t$: CIR is not Gaussian (unlike OU or
ABM), since $\sqrt{X_t}$ scales the noise unevenly across the state
space. Recovering the full law requires the noncentral chi-squared
transition density flagged above, not this ODE method.
\end{remark}

%--------------------------------------------------------------
\section{The Feller condition}
%--------------------------------------------------------------

\begin{theorem}[Feller condition]\label{thm:feller}
If $2\kappa\theta \ge \sigma^2$, then $X_t>0$ for all $t\ge0$, almost
surely, given $X_0>0$. If $2\kappa\theta < \sigma^2$, $X_t$ can reach
$0$ in finite time (though it is instantaneously reflected back into
$(0,\infty)$ by the drift, since $\drift(t,0)=\kappa\theta>0$).
\end{theorem}

\begin{proof}[Citation]
This is Feller's boundary classification for one-dimensional
diffusions, applied to \eqref{eq:cirsde}; see \citeasnoun{cir1985} or
\citeasnoun{brigomercurio2006}. \emph{(Cited, not proved --- the proof
requires the scale-function/speed-measure machinery for
one-dimensional diffusions, which nothing else in this note depends
on.)}
\end{proof}

\begin{remark}[Why this condition matters here specifically]
Every other assumption in this companion (Lipschitz coefficients,
$\kappa>0$ for OU's stationary law) has been a technical convenience.
This one is load-bearing in a different way: \eqref{eq:cirsde} is not
even well-defined pathwise if $X_t$ goes negative, since
$\sqrt{X_t}$ has no real value there. The Feller condition is what
guarantees the SDE never has to confront that question in continuous
time. Nothing in \texttt{hasklib} currently enforces
$2\kappa\theta\ge\sigma^2$ on constructing a \texttt{CIR} instance;
whether to validate it, and what to do if a calibrated parameter set
violates it (a common occurrence in practice), is a live design
question, not a settled one.
\end{remark}

\begin{remark}[Relevance beyond the pricing showcase]
CIR is, alongside Hull--White, one of the classical short-rate models
in \citeasnoun{brigomercurio2006} --- valued for keeping rates
strictly positive when the Feller condition holds, in contrast to
Hull--White's Gaussian (and hence potentially negative) short rate.
The parameter tension is real: market-calibrated CIR parameters
frequently fail the Feller condition, which is precisely why the
numerical scheme discussed next is not a mere implementation detail.
\end{remark}

%--------------------------------------------------------------
\section{Motivation for full truncation}
%--------------------------------------------------------------

Even when the Feller condition holds and the true process $X_t$
never reaches zero, a naive Euler--Maruyama discretisation
(\texttt{schemes.tex}) can still fail: the update
\[
   \hat X_{t+\Delta t} = \hat X_t + \kappa(\theta-\hat X_t)\,\Delta t
                          + \sigma\sqrt{\hat X_t}\,\sqrt{\Delta t}\,Z
\]
has no mechanism preventing $\hat X_{t+\Delta t}<0$ for an unlucky
draw of $Z$ --- and once that happens, $\sqrt{\hat X_{t+\Delta t}}$ at
the \emph{next} step is not a real number at all. This is a
discretisation artefact: it can occur even when the Feller condition
guarantees the true continuous-time process stays positive, simply
because a discrete-time approximation is not required to inherit that
guarantee.

\begin{definition}[Full truncation scheme]\label{def:fulltrunc}
Write $x^+=\max(x,0)$. The full truncation Euler scheme
\citeasnoun{lord2010} replaces \emph{every} appearance of $\hat X_t$
inside the drift and diffusion evaluations --- but not the running
value being updated itself --- by its positive part:
\[
   \hat X_{t+\Delta t} = \hat X_t + \kappa\big(\theta-\hat X_t^+\big)\,\Delta t
                          + \sigma\sqrt{\hat X_t^+}\,\sqrt{\Delta t}\,Z .
\]
\end{definition}

\begin{remark}[Why \emph{both} terms are truncated, not just the diffusion]
An earlier, more obvious fix (\emph{partial truncation}) replaces
$\hat X_t$ by $\hat X_t^+$ only inside the square root, leaving the
drift's mean-reversion term $\kappa(\theta-\hat X_t)$ using the raw
(possibly negative) value. \citeasnoun{lord2010} found this still
carries a systematic bias, and their contribution was to show that
truncating the drift's evaluation \emph{as well} --- \Cref{def:fulltrunc}
--- reduces that bias further, at no extra implementation cost. The
running value $\hat X_t$ itself is still allowed to go negative and
stay negative for a stretch (it is only clipped where it is
\emph{used}, inside the coefficient evaluations); this is a
deliberate feature of the scheme, not an oversight, and is what
distinguishes ``full'' from ``partial'' truncation.
\end{remark}

\begin{remark}[Correspondence with the code]
This is the discretisation \texttt{hasklib}'s roadmap specifies for
CIR (OU~$\to$~ABM~$\to$~CIR with full truncation~$\to$~Milstein), in
preference to plain \texttt{EulerMaruyama::step} applied naively:
wherever \texttt{CIR::drift}/\texttt{CIR::diffusion} are evaluated
inside a stepping scheme, the state passed in should be $\hat X_t^+$,
not $\hat X_t$ directly. Establishing the strong-order convergence
properties of this scheme is a separate question from motivating it
and is not addressed here.
\end{remark}

%--------------------------------------------------------------
\bibliographystyle{agsm}
\bibliography{refs}

\end{document}